\documentclass[12pt,a4paper]{letter}
\usepackage[utf8]{inputenc}
\usepackage[T1]{fontenc}
\usepackage{lmodern}
\usepackage[margin=1in]{geometry}

\signature{Ian Todd\\Sydney Medical School\\University of Sydney}
\address{Ian Todd\\Sydney Medical School\\University of Sydney\\Sydney, NSW, Australia\\itod2305@uni.sydney.edu.au}

\begin{document}

\begin{letter}{Editor\\New Journal of Physics}

\opening{Dear Editor,}

I am writing to submit the manuscript ``Information Geometry of the Quantum-Classical Transition: From Photosynthetic Excitons to Neural Oscillations'' for consideration in \textit{New Journal of Physics}.

This paper connects two areas well represented in NJP---quantum information geometry and quantum biology---in a way that, to our knowledge, has not been done before. We prove that decoherence is a dimensional collapse in the quantum Fisher information manifold: the dephasing channel maps a $(d^2-1)$-dimensional quantum state manifold onto the $(d-1)$-dimensional classical probability simplex, and we characterize the geometry of this collapse exactly for qubits (closed-form separation angle) and completely for $d \geq 3$ (global orthogonality theorem: the diagonal and off-diagonal tangent subspaces are Bures-orthogonal \emph{only} at diagonal states---any coherence breaks orthogonality).

The central application is to photosynthetic energy transfer. NJP published the foundational papers on environment-assisted quantum transport (ENAQT)---Rebentrost et al.\ (2009, NJP 11:033003) and Plenio \& Huelga (2008, NJP 10:113019)---which showed that intermediate dephasing enhances transport efficiency. We apply the principal angle framework to two structurally distinct complexes: the Fenna--Matthews--Olson (FMO) complex ($d = 7$, collapse $48 \to 6$) and the PE545 cryptophyte antenna ($d = 8$, collapse $63 \to 7$). Using Haken--Strobl dynamics with trapping, we find that at the ENAQT-optimal dephasing rate, both complexes are geometrically close to classical---all principal angles exceed $87^\circ$ for FMO and $89^\circ$ for PE545---with the strong quantum-classical mixing regime lying at much weaker dephasing. The consistency of this geometric signature across complexes with different chromophore types, symmetries, and energy landscapes suggests that near-classical Bures geometry at the transport optimum is a generic feature of environment-assisted transport.

We extend the framework beyond quantum biology to neural oscillations. A cortical gamma-oscillation assembly ($d = 10$) sits ${\sim}12$ orders of magnitude past the quantum-classical boundary: at biological dephasing ($kT/\hbar \approx 4 \times 10^{13}$\,Hz), all principal angles are $90.00^\circ$ to machine precision. The framework quantifies the \emph{detection threshold}---the microenvironment shielding factor needed for any residual quantum structure to become geometrically distinguishable---converting the qualitative claim ``neural systems are classical'' into a quantitative geometric bound. This complements Tegmark's (2000) timescale argument with a continuous geometric measure.

The paper fits NJP's scope in ``biological and medical physics'' and complements the journal's existing contributions to quantum biology and quantum information geometry. The cross-scale application---from photosynthetic excitons to cortical oscillations---addresses NJP's scope in biological physics more broadly. We believe it would also be of interest to NJP's readership in quantum Darwinism and decoherence theory, given Zurek's NJP contributions to that program.

The manuscript has not been submitted elsewhere and is not under consideration at any other journal. There are no conflicts of interest.

\closing{Sincerely,}

\end{letter}
\end{document}
